%% kapa: exam=1 round=21 year=2010 period=1 form=A subject=economics subject_ko=경제원론 seq=20 number=60
%% source: Q-Net 감정평가사 기출문제 (cst00309) artlSeq=5024181 file=extracted/1008/5024181/2010년 제21회 감정평가사 1차시험 기출문제/1교시.hwp
%% conversion: kordoc 4.24.51 -> tools/gen1.py (auto); review=proofread-sonnet; flags=none
\begin{question}{60}
A기업의 총비용함수는 TC = 50 + 10Q이다. 이에 관한 설명으로 옳지 않은 것은? (단, TC는 총비용이고, Q는 생산량이다.)
\begin{choices}
  \item 생산량이 5일 때 평균고정비용은 10이다.
  \item 생산량이 5일 때 평균가변비용은 10이다.
  \item 생산량이 5일 때 한계비용은 10이다.
  \item 생산량이 증가할수록 한계비용이 상승한다.
  \item 생산량이 증가할수록 평균총비용이 하락한다.
\end{choices}
\end{question}
\answer{4}
